\documentclass[11pt,UTF8,fontset=fandol]{ctexart}
\usepackage{amsmath,amssymb,amsfonts}
\usepackage{graphicx,float,caption,subcaption}
\usepackage{booktabs,array,multirow,makecell}
\usepackage{geometry,enumitem,xcolor}
\usepackage[hidelinks]{hyperref}
\usepackage{url}
\geometry{margin=2.3cm}
\graphicspath{{/Users/yye/project/论文/assets/whendollm/}}
\captionsetup{font=small,labelfont=bf}
\setlength{\parskip}{0.4em}
\hypersetup{pdftitle={何时 LLM 能帮上节点分类 中文精读导读},pdfauthor={Wu et al.}}

% 小白知识卡片：上下加线 + 粗体标签，稳健不挑内容
\newenvironment{tip}[1]{%
  \par\smallskip\noindent\rule{\linewidth}{0.4pt}\par\nobreak
  \noindent\textbf{【小白补丁】#1}\par\nobreak\noindent\ignorespaces
}{%
  \par\noindent\rule{\linewidth}{0.4pt}\par\smallskip
}

\title{何时该让大模型来做图节点分类？一篇系统实证分析的中文精读导读}
\author{原文：Wu, Shen, Ge, Shan, Jiao, Sun, Cheng（ICML 2025）}
\date{}

\begin{document}
\maketitle
\begin{center}\small\textit{本导读为面向初学者的中文学习笔记：含要点翻译、框架图解读与关键结论数据重述。}\end{center}

\section*{论文卡片}
\begin{center}
\begin{tabular}{@{}ll@{}}
\toprule
\textbf{标题} & When Do LLMs Help With Node Classification? \\
            & A Comprehensive Analysis \\
\textbf{作者} & Xixi Wu, Yifei Shen, Fangzhou Ge, Caihua Shan, \\
            & Yizhu Jiao, Xiangguo Sun, Hong Cheng \\
\textbf{机构} & 香港中文大学、微软亚洲研究院、伊利诺伊大学厄巴纳-香槟分校 \\
\textbf{会议} & ICML 2025（国际机器学习大会） \\
\textbf{arXiv} & 2502.00829 \\
\textbf{代码} & \url{https://llmnodebed.github.io/} \\
\textbf{方向} & 图（Graph）$\times$ 大模型 $\times$ 系统性实证评测（Benchmark） \\
\textbf{一句话} & 不提新方法，而是把各种“用 LLM 做图节点分类”的套路拉到\\
              & 统一框架里公平比一遍，告诉你\emph{什么时候、用哪招}才真有用。 \\
\bottomrule
\end{tabular}
\end{center}

\section{一、这篇论文要回答什么问题}

最近两年，大家都在尝试用\textbf{大语言模型（LLM，Large Language Model，如 GPT-4o、LLaMA）}
去做一件传统上由\textbf{图神经网络（GNN）}负责的事情——\textbf{节点分类}（给图里的每个点打标签，
比如判断一篇论文属于哪个学科、一个用户是不是水军）。论文很多，效果看起来都不错，
但作者一针见血地指出一个尴尬现状：\textbf{缺乏清晰的设计指南}。

具体来说，工程师在真实项目里会犯难：
\begin{enumerate}[leftmargin=2em]
  \item LLM 又慢又贵（训练、推理成本远高于传统 GNN）。到底\textbf{什么情况下}多花这份钱才值得？
  \item “用 LLM 做节点分类”有好多\textbf{流派}（后面会讲到 Encoder/Explainer/Predictor 等），
        每一派各自适合什么场景？
  \item 已有的对比工作（如 GLBench）因为\textbf{骨干模型不统一、代码库不统一}，
        比出来的结论\textbf{不公平、不可靠}，没法直接拿来指导选型。
\end{enumerate}

\textbf{本文目标}：不发明新方法，而是做一次\textbf{公平、系统、可复现}的大比武。为此作者
（1）造了一个统一的测试平台 \textbf{LLMNodeBed}；（2）训练并评测了\textbf{超过 2{,}700 个模型}，
最终提炼出 \textbf{8 条可操作的经验（Takeaway）}，直接回答“何时该上 LLM、上哪一派”。
这是一篇典型的\textbf{分析类/基准类（benchmark）工作}，对工程选型极有参考价值。

\section{二、摘要中译}

\begin{quote}
节点分类是图分析中的一项基础任务，在众多领域有广泛应用。近来大语言模型（LLM）的突破，
催生了一批基于 LLM 的节点分类方法。尽管许多研究都展示了 LLM 方法令人瞩目的性能，
但\textbf{缺乏清晰的设计指南}可能会阻碍它们在实践中的落地。本文旨在通过对这些算法进行
\textbf{公平、系统的对比}来建立这样的指南。作为第一步，我们开发了 \textbf{LLMNodeBed}——
一个用于“基于 LLM 的节点分类”的综合代码库与测试平台。它包含 10 个同质图（homophilic）数据集、
4 个异质图（heterophilic）数据集、8 种基于 LLM 的算法、8 个经典基线，以及 3 种学习范式。
随后，我们开展了大规模实验，\textbf{训练并评测了超过 2{,}700 个模型}，以确定影响性能的关键\emph{设定}
（如学习范式、同质性）与关键\emph{组件}（如模型规模、提示词）。我们的发现总结出 \textbf{8 条洞见}，例如：
（1）在\textbf{半监督}设定下，基于 LLM 的方法能显著超过传统方法，但在\textbf{监督}设定下这种优势变得微弱；
（2）在\textbf{零样本}设定下，图基础模型（GFM）能击败开源 LLM，但仍不及 GPT-4o 这类强 LLM。
我们希望 LLMNodeBed 的发布连同这些洞见，能推动可复现研究并启发该领域的后续工作。
代码与数据集发布于 \url{https://llmnodebed.github.io/}。
\end{quote}

\section{三、核心贡献与三句话结论}

\textbf{三大贡献}：
\begin{enumerate}[leftmargin=2em]
  \item \textbf{一个测试平台}：发布 LLMNodeBed（基于 PyG），统一了骨干模型与代码，
        初版含 14 个数据集、8 种 LLM 算法、3 种学习范式，一条命令即可跑全部实验、自动生成论文里所有表格。
  \item \textbf{一次大规模公平实验}：训练评测 $>2{,}700$ 个模型，系统分析了\emph{学习范式、同质性、
        语言模型类型与规模、提示词设计}如何影响各类算法。
  \item \textbf{8 条洞见与实操建议}：明确指出每一类算法\emph{在什么设定下最强}、关键组件怎么选，
        给出直观解释与“避坑”提示。
\end{enumerate}

\textbf{如果只记三句话}：
\begin{enumerate}[leftmargin=2em]
  \item \textbf{标签少（半监督）时，LLM 真香}；\textbf{标签多（监督）时，老老实实用 GNN 就够了}，
        LLM 的提升变得很小、未必划算。
  \item \textbf{图“没什么用”时（异质图、结构信息弱），LLM 当编码器优势巨大}；
        图本身信息很强时，换不换 LLM 差别不大。
  \item \textbf{零样本时，把图结构喂给强 LLM（如 GPT-4o）直接推理效果最好}；专门训练的图基础模型还追不上它。
\end{enumerate}

\section{四、预备知识（小白先看这里）}

\begin{tip}{什么是节点分类（Node Classification）}
一张图由\textbf{节点}（点，代表实体，如论文、用户、商品）和\textbf{边}（连线，代表关系，如引用、关注、共同购买）
组成。节点分类就是：已知一部分节点的标签，去\textbf{预测其余节点的标签}。本文研究的是
\textbf{文本属性图（TAG，Text-Attributed Graph）}——每个节点还自带一段文字（如论文的标题摘要、
商品的描述），正是这段文字给了 LLM 用武之地。
\end{tip}

\begin{tip}{三种“用 LLM”的范式：Encoder / Explainer / Predictor}
这是全文最重要的分类，务必记住（对应图~\ref{fig:overview}）：
\begin{itemize}[leftmargin=1.4em]
  \item \textbf{LLM-as-Encoder（当编码器）}：让 LLM 把每个节点的文字编码成\textbf{更好的特征向量}，
        再喂给一个 GNN 去分类。LLM 只负责“把文字变成好用的数字”。代表方法 ENGINE。
  \item \textbf{LLM-as-Explainer（当解释器/增强器）}：让 LLM 用生成能力\textbf{给节点文字“扩写、补充解释”}
        （例如先让它预测+给理由），把原文本变丰富，再走“编码 + GNN”。代表方法 TAPE。
  \item \textbf{LLM-as-Predictor（当预测器）}：把图结构信息也塞进 LLM，让 LLM \textbf{直接吐出标签}。
        代表方法 GraphGPT、LLaGA，以及纯指令微调 LLM（本文记作 $\mathrm{LLM_{IT}}$）。
\end{itemize}
通俗记忆：Encoder“只出特征”、Explainer“先扩写文本再出特征”、Predictor“一条龙直接出答案”。
\end{tip}

\begin{tip}{什么是同质性（Homophily）——本文的关键变量}
同质性 = “物以类聚”的程度，指\textbf{一条边两端的节点有多大概率属于同一类}。
\textbf{同质图（homophilic）}：相连的点大多同类（如论文引用网络，引用关系常发生在同领域论文间，
本文中 Cora 同质率 82.5\%）。此时\textbf{图结构本身就是强信号}，GNN 沿着边传播就能分对很多。
\textbf{异质图（heterophilic）}：相连的点常常不同类（本文中 Texas 等网页数据集同质率仅 6.7\%）。
此时图结构帮不上忙，甚至会误导，模型更要依赖\textbf{节点自身的文字内容}——这正是 LLM 编码器的强项。
记住一句话：\textbf{图越“没用”（同质性越低），LLM 越能帮上忙}。
\end{tip}

\begin{tip}{三种学习设定与零样本里的两条路}
\textbf{半监督（Semi-supervised）}：只给\emph{极少}带标签的节点（如每类几个），考验“少标签也能学好”。\\
\textbf{监督（Supervised）}：给\emph{大量}带标签节点（本文多数数据集用 60\% 训练）。\\
\textbf{零样本（Zero-shot）}：\emph{完全不给}目标图的标签去训练。零样本下又分两条路：
（a）\textbf{Direct Inference 直接推理}——精心设计提示词（Prompt）让 LLM 当场分类；
（b）\textbf{GFM（图基础模型，Graph Foundation Model）}——先在别的大图上预训练好，再迁移到目标图。
\end{tip}

形式化记号（看不懂可跳过）：文本属性图记作 $G=(V,E,S)$，其中 $V$ 是节点集、$E$ 是边集、
$S$ 是每个节点 $v$ 携带的文字。有标签节点记 $V_l$、无标签节点记 $V_u$。
传统做法把文字用 bag-of-words 或 TF-IDF 压成“浅层嵌入” $X\in\mathbb{R}^{|V|\times d}$ 喂给 GNN；
而 LLM 方法直接吃\textbf{原始文字}，指望预训练知识带来增益。

\section{五、实验设计与数据集（怎么做到“公平”）}

作者用图~\ref{fig:overview} 这张总览图把整个测试平台 LLMNodeBed 讲清楚了。

\begin{figure}[H]\centering
  \includegraphics[width=0.95\linewidth]{cand_p02_k00.pdf}
  \caption{LLMNodeBed 总览（原文 Figure 1，图内英文为原文保留）。自上而下分五层：
  \textbf{Datasets} 数据集——覆盖学术、网页、社交、电商、异质图五大领域；
  \textbf{Backbones} 骨干模型——GNN（GCN/SAGE/GAT/H2GCN）、小型语言模型 LM（SenBERT/RoBERTa）、
  大模型 LLM（Qwen/Mistral/LLaMA）；
  \textbf{Supported Methods} 支持的方法——半监督/监督下分 \emph{LLM-as-Encoder、LLM-as-Explainer、
  LLM-as-Predictor} 三派，零样本下分\emph{带提示词的直接推理}与\emph{图基础模型 GFM}；
  \textbf{Evaluation} 评测——准确率 Accuracy、Macro-F1、幻觉率 Hallucination。
  统一骨干 + 统一代码，是“公平对比”的基础。}
  \label{fig:overview}
\end{figure}

\textbf{“公平”的关键设计}：在主对比里，所有方法都用\textbf{同一套组件}——GNN 一律用 GCN，
小模型 LM 一律用 RoBERTa-355M，大模型 LLM 一律用 Mistral-7B。这样性能差异就只来自\emph{方法设计本身}，
而不是“谁用了更强的骨干”。每个实验跑 4 次取平均。

\textbf{数据集（源自原文 Table 1，节选关键列）}：作者特意挑了同质性差异巨大的数据集，
好观察“同质性”这个变量的影响。

\begin{table}[H]\centering
\caption{部分数据集统计（数值源自原文 Table 1）。同质率越低，图结构越“不可靠”。}
\begin{tabular}{lcccccc}
\toprule
 & \textbf{Cora} & \textbf{Pubmed} & \textbf{WikiCS} & \textbf{Reddit} & \textbf{Computer} & \textbf{Texas} \\
\midrule
领域 & 学术 & 学术 & 网页 & 社交 & 电商 & 网页(异质) \\
节点数 & 2{,}708 & 19{,}717 & 11{,}701 & 33{,}434 & 87{,}229 & 187 \\
同质率(\%) & 82.52 & 79.24 & 68.67 & 55.52 & 85.28 & \textbf{6.69} \\
\bottomrule
\end{tabular}
\end{table}

\section{六、统一分析框架：用“互信息”解释为什么}

这篇论文虽然是实证分析，但它给出了一个\textbf{贯穿全文的理论视角}来解释各种现象——
\textbf{互信息（mutual information）分解}。这是理解全文结论的“总钥匙”，值得专门讲一节。

节点分类（用交叉熵损失训练）本质上是在\textbf{最大化“图 $G$”与“标签 $Y_l$”之间的互信息} $I(G;Y_l)$。
作者把图 $G$ 看成“节点属性 $X$（文字）”和“结构 $E$（边）”的联合分布，于是这个互信息可以拆成两项：
\begin{equation*}
  I(G;Y_l)=I(X,E;Y_l)=\underbrace{I(E;Y_l)}_{\text{结构提供的信息}}
  +\underbrace{I(X;Y_l\mid E)}_{\text{文字额外提供的信息}}. \tag{1}
\end{equation*}

\begin{tip}{这个公式怎么读、有什么用}
\textbf{互信息} $I(A;B)$ 衡量“知道 $A$ 能帮你确定 $B$ 多少”，值越大越相关。公式 (1) 把“做对分类需要的总信息”
拆成两块：
\begin{itemize}[leftmargin=1.4em]
  \item 第一项 $I(E;Y_l)$：\textbf{光靠图的结构（边）能提供多少线索}——这是传统 GNN 吃的饭。
        同质图里这一项很大（沿边传播就够）。
  \item 第二项 $I(X;Y_l\mid E)$：\textbf{在已知结构之后，节点文字还能额外补多少线索}——这是 LLM 吃的饭。
\end{itemize}
\textbf{一句话推论}：当第一项（结构）已经很大时，第二项（文字）能发挥的空间就小，LLM 帮不上太多；
当第一项很小时（如异质图、或标签很少导致学不到结构规律），第二项就成了主力，\textbf{LLM 价值凸显}。
后面几乎所有结论，都能用这把钥匙解释。
\end{tip}

举两个例子帮助体会：
\begin{itemize}[leftmargin=2em]
  \item \textbf{为什么半监督下 LLM 更有用}？标签极少时，模型很难从稀疏标签里学到可靠的“结构规律”，
        即 $I(E;Y_l)$ 实际可利用得很低，于是文字项 $I(X;Y_l\mid E)$ 成了救命稻草——LLM 的语义理解大显身手。
  \item \textbf{为什么监督下强 LLM 优势缩小}？标签充足时 $I(E;Y_l)$ 被充分利用、变得很大，
        对文字的依赖下降，于是“用 GPT-4o 还是用小一点的模型”差别就不明显了。
\end{itemize}

\section{七、关键发现与结果表解读（8 条 Takeaway）}

这是本文精华。作者把 8 条洞见嵌在实验里，下面按“何时有用 / 何时无用”重新组织，并配关键数字。

\subsection{第一组：半监督 vs.\ 监督——LLM 到底什么时候值}

\textbf{主结果（数值源自原文 Table 2，准确率 \%，节选关键方法 + 平均分 Avg.）}：

\begin{table}[H]\centering
\caption{半监督与监督设定下的性能（源自原文 Table 2，准确率 \%）。Avg.\ 为多数据集平均；越高越好。}
\begin{tabular}{llcc}
\toprule
设定 & 代表方法（类别） & 半监督 Avg. & 监督 Avg. \\
\midrule
经典 & GAT（纯 GNN，浅层嵌入） & 74.67 & 78.70 \\
经典 & RoBERTa（纯文本 LM） & 68.06 & 78.17 \\
\midrule
LLM-as-Encoder & GCN$_{\mathrm{LLMEmb}}$ & 78.51 & 81.09 \\
LLM-as-Encoder & ENGINE & 78.40 & 81.32 \\
LLM-as-Explainer & TAPE & \textbf{78.43} & \textbf{81.80} \\
LLM-as-Predictor & LLaGA & 75.07 & 81.69 \\
LLM-as-Predictor & $\mathrm{LLM_{IT}}$（纯微调 LLM） & 65.76 & 72.38 \\
\bottomrule
\end{tabular}
\end{table}

\textbf{怎么读这张表（对应 Takeaway 1{--}5）}：
\begin{itemize}[leftmargin=2em]
  \item \textbf{T1：合理引入 LLM 总能提升}。最好成绩几乎都来自 LLM 方法，且这一结论在百万级大图
        （ogbn-products）上同样成立——用 LLM 挖掘文字信息确实有用。
  \item \textbf{T2：半监督的提升远大于监督}。从经典 GAT 到 LLM 方法，半监督下 Avg.\ 从 74.7 涨到约 78.5
        （提升明显）；而监督下大家都挤在 81 至 82 分，\textbf{LLM 相对经典 GNN 的优势变得微弱}。
        （这正是公式 (1) 的推论。）
  \item \textbf{T3：标签强依赖文字时，Explainer（TAPE）最猛}。在 Cora、WikiCS 等“结构没那么重要”的
        学术/网页数据上 TAPE 拿第一/第二；但在需要深层结构理解的社交网络（如 Reddit 上判断高热度用户）就吃力。
  \item \textbf{T4：Encoder 是“性价比之王”}。LLM-as-Encoder 在所有数据集上都表现稳健，
        且耗时远小于要“生成长解释”的 Explainer，在\textbf{效率与精度间取得平衡}。
  \item \textbf{T5：Predictor 要靠“标签多”才行}。监督下 LLaGA 很强（10 个数据集里 5 个最佳），
        但半监督下 Predictor 不稳定（Macro-F1 低、输出类别分布失衡）。尤其\textbf{纯微调 LLM（$\mathrm{LLM_{IT}}$）
        最弱且最慢}（arXiv 上一轮要 30+ 小时），说明\textbf{光有 LLM 不行，必须把图结构喂进去}。
\end{itemize}

\subsection{第二组：零样本——不给标签时谁最强}

\textbf{零样本结果（数值源自原文 Table 3，准确率 \%，4 数据集平均 Avg.）}：

\begin{table}[H]\centering
\caption{零样本设定下的准确率（源自原文 Table 3，\%，4 数据集平均）。}
\begin{tabular}{llc}
\toprule
类型 & 方法 & Avg.\ Acc \\
\midrule
强 LLM (GPT-4o) & Direct 直接提示 & 61.30 \\
强 LLM (GPT-4o) & \textbf{w.\ Summary（先总结邻居）} & \textbf{64.99} \\
开源 LLM (LLaMA-8B) & Direct & 49.56 \\
开源 LLM (LLaMA-8B) & w.\ Summary & 53.27 \\
GFM 图基础模型 & ZeroG & 55.56 \\
GFM (迁移的 Predictor) & LLaGA & 28.53 \\
\bottomrule
\end{tabular}
\end{table}

\textbf{怎么读（对应 Takeaway 6{--}7）}：
\begin{itemize}[leftmargin=2em]
  \item \textbf{T6：GFM 能赢开源 LLM，但赢不了 GPT-4o}。ZeroG（55.6）大体优于 LLaMA-8B（约 50），
        却明显不及 GPT-4o（约 61 至 65）。而把训练好的 Predictor（LLaGA）拿去跨数据集迁移时\textbf{彻底崩盘}
        （仅 28.5），因为它的“投影器”是\emph{数据集专用}的，换个图就失灵。
  \item \textbf{T7：给 LLM 喂结构信息能涨分，但“花式提示词”几乎没用}。CoT、ToT、ReAct 这些
        \emph{通用推理}提示词对节点分类提升很小（LLaMA 遇到复杂提示甚至会输出超长、难解析的内容）；
        真正有效的是\textbf{把邻居信息塞进去}——“w.\ Neighbor”尤其是\textbf{“w.\ Summary”（先让 LLM 总结邻居再分类）}，
        在同质图（Cora、Photo）上能涨 3\% 到 10\%。原因还是同质性：邻居和自己同类，邻居信息就是强线索。
\end{itemize}

\subsection{第三组：每一派内部的精细分析}

\textbf{T8：图越“没用”，LLM 当编码器优势越大（源自原文 Table 4{--}5）}。
对比“用 LLM 编码 vs.\ 用小模型 LM 编码”，在同质图上二者差距很小；但在\textbf{异质图}上，
LLM 编码器优势巨大——例如 Texas 数据集上，配合 GCN/SAGE/H2GCN，LLM 编码比 LM 编码\textbf{准确率高出约 12\%}，
跨方法差距可达 10\%。再次印证：结构帮不上忙时（异质图），就得靠文字，而 LLM 最会读文字。

\textbf{Explainer 用强 LLM 划算吗？（Table 6{--}7）}：把 TAPE 里的 Mistral-7B 换成 GPT-4o，
半监督下确有提升（Pubmed 上 +4\%），但\textbf{监督下提升微乎其微}。
\textbf{建议：标签充足时，用便宜的开源 LLM 当 Explainer 即可}，没必要为 GPT-4o 多花钱。

\textbf{Predictor 选多大的模型？（Table 8{--}9）}：在 LLaGA 上把 Qwen 从 3B 加到 32B，
性能一开始随规模上升，但\textbf{过了 7B/14B 就基本饱和}，而 32B 推理慢 5 倍（每样本 200+ 毫秒）。
\textbf{建议：Predictor 选 Qwen-7B 或 14B 最划算；同规模里 Mistral-7B 综合最优}，是推荐骨干。

\subsection{别忽视的两个工程现实：成本与幻觉}

\begin{tip}{成本（源自原文成本分析）：LLM 方法普遍又慢又费显存}
经典方法极快：GNN 几秒收敛、LM 微调几分钟，推理可在毫秒级处理上千样本，适合工业实时部署。
而 LLM 方法：Explainer 最慢（TAPE 在 Pubmed 上要 5.9 小时，因要给每个节点生成解释文本）；
Encoder 在 LLM 方法里最快（十几分钟）且显存与普通 GNN 相当；Predictor 要数小时训练。
\textbf{推理时 LLM 方法与经典方法仍有巨大效率鸿沟}——这正是选型时必须权衡的“性价比”。
\end{tip}

\begin{tip}{幻觉（hallucination）：Predictor 会“答非所问”}
让 LLM 直接当 Predictor 时，它可能输出\textbf{标签集合里根本不存在的类别}（即幻觉）。
论文附录专门统计了各 Predictor 方法的平均幻觉率，并展示了 GraphGPT 在 Citeseer/Pubmed 上的幻觉例子。
这是“让 LLM 直接出标签”这条路独有的风险，Encoder/Explainer 路线因为最后用的是 GNN 分类头，不存在此问题。
\end{tip}

\section{八、关键结论速查表（何时有用 / 何时无用）}

把全文 8 条洞见浓缩成一张“选型速查表”：

\begin{table}[H]\centering
\caption{LLM 节点分类“何时用、用哪招”速查（本导读据原文 8 条 Takeaway 重排）。}
\renewcommand{\arraystretch}{1.25}
\begin{tabular}{p{4.2cm}p{4.0cm}p{4.2cm}}
\toprule
\textbf{你的场景} & \textbf{推荐做法} & \textbf{为什么 / 注意} \\
\midrule
标签很少（半监督） & 上 LLM，优先 \textbf{Encoder} 或 \textbf{Explainer(TAPE)} & 提升最明显；文字信息此时最值钱 \\
标签很多（监督） & \textbf{经典 GNN 通常就够} & LLM 优势微弱，未必抵得过成本 \\
异质图 / 结构没用 & \textbf{LLM-as-Encoder} & 比小模型 LM 高出可达 12\%，优势最大 \\
要均衡精度与成本 & \textbf{LLM-as-Encoder} & LLM 里最快、显存最省，性价比之王 \\
标签足 + 想直接出标签 & \textbf{LLaGA}（Predictor） & 强但要喂图结构；纯微调 LLM 最差最慢 \\
零样本、没法训练 & \textbf{强 LLM(GPT-4o) + 邻居/总结提示} & 比 GFM 更好；花式推理提示几乎无用 \\
零样本且想用 GFM & \textbf{ZeroG} & 胜开源 LLM、不及 GPT-4o；迁移型 Predictor 会崩 \\
Explainer 选骨干 & 标签足时用\textbf{便宜开源 LLM} 即可 & 换 GPT-4o 在监督下几乎不涨分 \\
Predictor 选规模 & \textbf{Qwen-7B/14B} 最划算 & 过了就饱和，32B 慢 5 倍 \\
\bottomrule
\end{tabular}
\end{table}

\section{九、小白知识点卡片（速查）}
\begin{itemize}[leftmargin=2em]
  \item \textbf{节点分类（Node Classification）}：给图里的点打标签的任务。
  \item \textbf{文本属性图 / TAG}：每个节点都带一段文字的图。
  \item \textbf{同质性（Homophily）}：相连节点同类的概率；高=结构有用，低=要靠文字。
  \item \textbf{LLM-as-Encoder}：让 LLM 把文字编码成好特征，再交给 GNN 分类。
  \item \textbf{LLM-as-Explainer}：让 LLM 先把节点文字“扩写补充”，增强后再分类（如 TAPE）。
  \item \textbf{LLM-as-Predictor}：把图喂进 LLM，让它直接输出标签（如 GraphGPT、LLaGA）。
  \item \textbf{GFM（图基础模型）}：在大图上预训练、可迁移到新图的模型（如 ZeroG）。
  \item \textbf{半监督 / 监督 / 零样本}：分别是“极少标签 / 大量标签 / 无标签训练”。
  \item \textbf{互信息 $I(\cdot;\cdot)$}：衡量两个量相关程度；本文用它解释“结构信息 vs.\ 文字信息”谁主导。
  \item \textbf{幻觉（Hallucination）}：模型输出不存在的类别/内容，是 Predictor 路线的风险。
  \item \textbf{Macro-F1}：对各类别一视同仁的 F1 平均，能暴露“类别不平衡时”的虚高准确率。
  \item \textbf{CoT / ToT / ReAct}：思维链/思维树/推理-行动等通用提示词技巧。
\end{itemize}

\section{十、一句话总结与谱系定位}

\textbf{一句话}：这篇论文不造新模型，而是搭了个统一擂台 LLMNodeBed、跑了 2{,}700+ 个模型，
用一把“互信息分解”的钥匙把结论讲透——\textbf{标签少 / 图没用（异质）时该上 LLM（尤其当 Encoder），
标签多时经典 GNN 就够；零样本时把图结构喂给强 LLM 直接推理最香}。

\textbf{在“图 $\times$ 大模型”谱系里的位置}：本系列里大多数论文是在\emph{提出一种新方法}——
G-Retriever 做图上 RAG、LLaGA 做投影对齐、GraphGPT/InstructGraph 做指令微调、G1 做强化学习……
而\textbf{本文是一篇“元层面”的分析/基准工作}：它把上述各类方法（Encoder/Explainer/Predictor/GFM）
\textbf{拉到同一标准下公平横评}，回答“到底该选谁”。因此它最适合作为\textbf{工程选型手册}与\textbf{入门地图}：
读完它，你再去看具体某个方法的论文，就知道那个方法属于哪一派、强在什么设定下。

\textbf{动手复现}：官方代码与数据 \url{https://llmnodebed.github.io/}（一条命令复现全部表格），
是系统了解“LLM 做图节点分类”全景的极佳起点。

\end{document}
